\documentclass{article}
\usepackage{amsmath}
\usepackage{algorithm}
\usepackage{algorithmic}

\begin{document}

\section{Cálculo de raíces}

\begin{equation}
f(x) = x \cdot \ln(2)
\label{eq:calc:f}
\end{equation}

\begin{equation}
a = -1
\label{eq:calc:a}
\end{equation}

\begin{equation}
b = 10
\label{eq:calc:b}
\end{equation}

\begin{algorithm}
\caption{Método de dicotomía}
\label{alg:calc:dicotomia}
\begin{algorithmic}
\STATE $tol = 1e-6$
\WHILE{$b - a > tol$}
    \STATE $c = (a + b)/2$
    \IF{$f(a) \cdot f(c) < 0$}
        \STATE $b = c$
    \ELSE
        \STATE $a = c$
    \ENDIF
\ENDWHILE
\RETURN $c$
\end{algorithmic}
\end{algorithm}

\begin{equation}
    c = 
    -0.000000
    \label{eq:res:c}
\end{equation}

\section{Cálculo de soluciones de $x^a + y^b = z^c$}

Valores $a$, $b$, y $c$
\begin{equation}
a = 1
\label{eq:calc:a}
\end{equation}

\begin{equation}
b = 2
\label{eq:calc:b}
\end{equation}

\begin{equation}
c = 2
\label{eq:calc:c}
\end{equation}

Valor máximo para las soluciones ($x, y, z \in \{1,2,..max\}$):
\begin{equation}
max = 10
\label{eq:calc:max}
\end{equation}

\begin{algorithm}
\caption{Número de soluciones}
\label{alg:calc:numsols}
\begin{algorithmic}
    \STATE sols = 0
    \STATE z = 1
    \WHILE{$z \leq max$}
        \STATE y = 1
        \WHILE{$y \leq max$}
            \STATE x = 1
            \WHILE{$x \leq max$}
                \IF{$x^a + y^b = z^c$}
                    \STATE sols = sols + 1
                    \STATE print("x =", x, " y =", y, "z=",z)
                \ENDIF
                \STATE x = x + 1
            \ENDWHILE
            \STATE y = y + 1
        \ENDWHILE
        \STATE z = z + 1
    \ENDWHILE
    \RETURN sols
\end{algorithmic}
\end{algorithm}

\begin{equation}
    sols = 
    5
    \label{eq:res:sols}
\end{equation}


\section{Ejemplo con ELSIFs}

\begin{algorithm}
\caption{Ejemplo de un algoritmo con ELSIFs anidados}
\label{alg:calc:Ejemplo con ELSIF}
\begin{algorithmic}
        \STATE $i = 10$
        \STATE $j = \sqrt{i^2 + 5}$
        \IF{$i \ge 5$}
            \STATE $i = i - 1$
            \WHILE{$j \le 20$}
                \STATE $j = j + \log(i)$
            \ENDWHILE
        \ELSIF{$i \le 3$}
            \STATE $i = i + 2$
        \ENDIF
        \RETURN $j$
\end{algorithmic}
\end{algorithm}



\begin{equation}
    j = 
    21.233074
    \label{eq:res:j}
\end{equation}

\section{Prueba con raíces cúbicas}
\begin{equation}
    a = \sqrt[3]{8}
    \label{eq:calc:a}
\end{equation}

\begin{equation}
    a = 
    2.000000
    \label{eq:res:a}
\end{equation}

\end{document}

